\section{Calibrate}

\textbf{Today.} Calibration is done by hand, as Section 3.5 described: starting from the results of an earlier run, the analyst checks whether production, consumption or both deviate from the reference values, decides whether the cause lies in the data or in the parameters, adjusts the parameter that seems least plausible, and runs the model again, over several days.
It has two parts.
The starting values come first: base prices and seasonal price spreads can be researched, while production costs cannot, because no data on them exist.
Then comes the loop of runs and adjustments, in which a new climate scenario in the 2023 data changed between about 1,700 and 2,050 calibrated values.

\textbf{What the step needs.} Each pass through the loop combines four kinds of work.
Computing how far each country's production and consumption are from their reference values is arithmetic.
Diagnosing why is reasoning: the documentation lists three possible causes when Russia exports too little, and six when China consumes too little in a future year, ranging from domestic production through willingness to pay and market power to infrastructure costs and expansion limits \citep[p.~26]{egging2019global}.
Choosing which parameter to change is judgment, guided by what is plausible and by the requirement that the choice should still make sense in another scenario.
Running the model is mechanical.
Because trade links all markets, a change for one country moves others, so the loop does not converge country by country.

\textbf{What an agent could do.} Calibration is the one step whose path cannot be written down in advance: which parameter to change next depends on what the last run showed.
In the terms of Section 4.1, it needs an agent rather than a workflow, but an agent working inside a harness of ordinary code.
The harness would compute a deviation report after every run, which the 2019 version of the model produced and the public Julia version does not (Section 3.4).
It would apply the agent's proposed changes, reject those outside set bounds and run the model on local hardware, exchanging only results with the agent, as Lukas Barner insisted.
The bounds can be taken partly from the documentation, which states that capacity slack should be at least a few percent and not more than 7 to 10\% \citep[p.~28]{egging2019global}, and that elasticities are not adjusted in calibration \citep[p.~26]{egging2019global}.
The agent's own task would be the diagnosis: read the signed deviations, reason about their likely cause, propose a change and write down why.
This is the design of the grey-box calibration study in Section 2.3, in which the LLM proposes a parameter vector with a rationale and a harness runs the model \citep{gomezguillen2026agentic}.
The hydrology benchmark in the same section shows which failures to guard against: stopping too early, proposing implausible values, and diagnosing a problem correctly but adjusting the wrong parameter \citep{li2026hydroagent}.
The first two can be prevented by code, through a stopping rule and bounds; the third can only be caught by reading the agent's reasoning, which is why every change must be logged with its rationale.
The starting values are a separate, simpler task: an agent could research reference prices and seasonal spreads, with sources, before the loop begins.
Here Lukas Barner pointed to a specific trap: in countries with subsidised domestic consumption, official gas prices are far below the willingness to pay the model needs, and LLMs are poor at telling the two apart.

What an agent would save is mainly the analyst's attention between runs.
A run of the 2023 model took 24 minutes on our laptop (Section 3.5), so a working day allows about twenty runs, and analysts today calibrate two scenarios in parallel to use the waiting time.
An agent could keep the loop going outside working hours, but more runs only help if each proposal is sensible; the gain in the grey-box study came from needing fewer model runs, not from faster ones.

\textbf{Where a person stays involved.} Three decisions stay with a person.
The first is to accept a calibration as good enough.
Reproducing the reference values does not show that the parameters behind the fit are right, just as matching an optimisation model's objective value does not prove its formulation correct \citep{li2026optargus}.
The second is to choose between fits of similar quality, which, as Section 2.6 argued, can embody different assumptions about the market and behave differently in a new scenario.
The third is market power, which the developers set by judgment and use to represent long-term contracts; an agent could flag where market power seems to drive a deviation, but should not change it.
In the terms of \citet{parasuraman2000model}, gathering, analysing and, within bounds, acting can be automated, while the decision to accept stays with a person.

The documentation also makes the agent testable, which is rare.
Because its diagnostic heuristics are written down, an agent's explanations of a deviation can be compared with them, rather than judged only by impression.
And the 2023 data allow a test with a known answer: perturb a calibrated scenario, let the agent restore the fit, and compare both the fit and the values it chose with the developers' own.
